\section{Main Code Components}
\label{sec:code}

\subsection{Code Base Architecture}
\label{sec:architecture}

All executable code lives under \texttt{source\_code/}. The report and this
chapter only describe that tree. Figure~\ref{fig:pipeline} is the runtime
path from the raw UCI file to the dashboard.

\begin{figure}[h]
  \centering
  \begin{tikzpicture}[
    node distance=3.6mm,
    box/.style={draw, rounded corners=2pt, fill=blue!5, align=center,
                font=\small, inner sep=3pt, minimum height=7.5mm},
    src/.style={draw, rounded corners=2pt, fill=orange!10, align=center,
                font=\small, inner sep=3pt, minimum height=7.5mm},
    arr/.style={-{Latex[length=1.8mm]}, thick}
  ]
    \node[src] (uci) {UCI Tetouan CSV};
    \node[box, below=of uci] (load) {\texttt{data\_loader.py}: parse and aggregate to hours};
    \node[box, below=of load] (feat) {\texttt{feature\_engineering.py}: calendar, lags, rolling, weather};
    \node[box, below=of feat] (model) {\texttt{models.py}: naive, ridge, forest, XGBoost, LightGBM};
    \node[box, below=of model] (eval) {\texttt{evaluation.py}: MAE, RMSE, WAPE, peak scores, figures};
    \node[box, below=of eval] (out) {\texttt{outputs/metrics.json} and \texttt{predictions.csv}};
    \node[box, below=of out] (app) {\texttt{app.py}: Streamlit forecast and threshold alert};
    \draw[arr] (uci) -- (load);
    \draw[arr] (load) -- (feat);
    \draw[arr] (feat) -- (model);
    \draw[arr] (model) -- (eval);
    \draw[arr] (eval) -- (out);
    \draw[arr] (out) -- (app);
  \end{tikzpicture}
  \caption{Pipeline from the public series to the alert dashboard.}
  \label{fig:pipeline}
\end{figure}

\begin{table}[h]
  \centering
  \small
  \begin{tabular}{@{}lp{9.4cm}@{}}
    \toprule
    Path & Responsibility \\
    \midrule
    \texttt{src/data\_loader.py} & Download the UCI file if it is not cached, standardise column names, sum the three zones, and aggregate to hourly demand. \\
    \texttt{src/feature\_engineering.py} & Build calendar flags, lags, shifted rolling statistics and temperature interactions. Cut the table into \(D_{\mathrm{train}}\), \(T_{\mathrm{valid}}\) and \(T_{\mathrm{test}}\). \\
    \texttt{src/models.py} & Fit the seasonal naive, ridge, random forest, XGBoost and LightGBM models. Choose hyperparameters on the validation block. \\
    \texttt{src/evaluation.py} & Compute MAE, RMSE, WAPE and the peak metrics, and write the figures listed in Sections~\ref{sec:eda} and~\ref{sec:results}. \\
    \texttt{run\_pipeline.py} & Run the four modules and store \texttt{outputs/metrics.json} plus \texttt{outputs/predictions.csv}. \\
    \texttt{app.py} & Streamlit page: test-block forecast, capacity line, breach-anticipation summary and an alert when the forecast exceeds \(\tau\). \\
    \bottomrule
  \end{tabular}
  \caption{Modules in \texttt{source\_code/}.}
  \label{tab:modules}
\end{table}

The dashboard reads saved predictions. It does not retrain. Changing a model
hyperparameter means re-running \texttt{run\_pipeline.py}, then refreshing
the Streamlit page.

\subsection{Tech Stack}
\label{sec:stack}

The language is Python 3.12. Table~\ref{tab:stack} lists the libraries pinned
in \texttt{source\_code/requirements.txt} and the role each one plays.

\begin{table}[h]
  \centering
  \small
  \begin{tabular}{@{}ll@{}}
    \toprule
    Library & Role in this project \\
    \midrule
    \texttt{pandas}, \texttt{numpy} & Tables, resampling, lags and rolling windows \\
    \texttt{scikit-learn} & Ridge, random forest, MAE during model selection \\
    \texttt{xgboost}, \texttt{lightgbm} & Gradient-boosting forecasters \\
    \texttt{matplotlib}, \texttt{seaborn} & EDA figures and the test-set charts \\
    \texttt{streamlit} & Forecast and alert dashboard \\
    \bottomrule
  \end{tabular}
  \caption{Python stack.}
  \label{tab:stack}
\end{table}

No extra service is required. The UCI file is fetched over HTTPS on the first
run and cached under \texttt{source\_code/data/raw/}. Later runs reuse the
cache.
